\chapter[Recommenders II: Deep Models]{Recommenders II: Deep Learning, Embeddings, and Retrieval}

\section{Learning objectives}
\begin{itemize}
  \item Explain why deep learning methods are useful for large-scale recommendation and retrieval problems in e-commerce.
  \item Distinguish candidate generation from re-ranking and understand why modern recommenders are often multi-stage systems.
  \item Describe embeddings, two-tower retrieval, sequence models, and session-aware recommendation at a system level.
  \item Analyze production concerns such as cold start, diversity, latency, and bias in deep recommender architectures.
\end{itemize}

\section{From classical recommenders to deep retrieval}
Chapter 7 focused on classical collaborative filtering, matrix factorization, and learning-to-rank.
Those methods remain important, but large-scale e-commerce platforms often face additional pressures:
\begin{itemize}
  \item extremely large catalogs,
  \item rapidly changing inventories,
  \item sparse user histories,
  \item multi-modal item information such as text and images,
  \item session behavior that changes within minutes,
  \item and tight latency requirements at serving time.
\end{itemize}

Deep learning is useful in this setting because it can learn richer representations from heterogeneous data.
Rather than relying only on manually designed similarity signals, deep systems can learn embeddings that organize users, items, queries, sessions, and contexts into a common representational space.

\section{Modern recommender architecture: retrieval and ranking}
Large recommender systems are rarely a single model.
They are usually multi-stage decision systems.

\subsection{Candidate generation}
Candidate generation or retrieval narrows a huge catalog to a much smaller set of promising items.
For a catalog containing millions of products, the system cannot fully rank every item in real time for every request.
Retrieval therefore acts as a coarse but scalable selection layer.

\subsection{Re-ranking}
After candidate generation, a second-stage model or rule system orders the retrieved candidates more precisely.
This stage can use richer features, business constraints, and context-aware logic because the item set is now small enough to handle with more expensive computation.

\subsection{Why multi-stage design matters}
The multi-stage architecture exists because of a systems tradeoff:
\begin{itemize}
  \item retrieval prioritizes speed and coverage,
  \item re-ranking prioritizes precision and business control.
\end{itemize}

Students should understand that recommendation quality depends not only on model accuracy, but also on how these stages interact.
A weak retrieval stage can permanently exclude relevant items before ranking ever sees them.

\section{Embeddings as learned representations}
Embeddings are dense vectors that represent entities such as users, items, queries, categories, or sessions.
They generalize the latent-factor intuition of matrix factorization, but in a more flexible framework.

\subsection{Why embeddings are powerful}
Embeddings are useful because they:
\begin{itemize}
  \item capture similarity in a continuous space,
  \item support nearest-neighbor retrieval,
  \item can combine multiple signals and modalities,
  \item transfer across tasks such as recommendation, search, and personalization.
\end{itemize}

\subsection{What can be embedded}
Common embedded entities in commerce include:
\begin{itemize}
  \item users or customer accounts,
  \item products,
  \item search queries,
  \item brands or categories,
  \item sellers,
  \item sessions,
  \item product text and images.
\end{itemize}

\subsection{Interpretation}
An embedding dimension does not usually have a simple human-readable meaning.
The value of the representation comes from geometry:
items or users that are close in vector space are more likely to be related in terms of behavior or semantics.

\section{Two-tower retrieval models}
One of the most influential deep retrieval architectures is the two-tower model.
It is widely used because it balances expressiveness with serving practicality.

\subsection{Basic structure}
A two-tower model learns:
\begin{itemize}
  \item one tower that maps user or context features to an embedding,
  \item one tower that maps item features to an embedding.
\end{itemize}

The compatibility between a user/context and an item is then computed using a similarity function such as dot product or cosine similarity.

\subsection{Why two towers are useful}
Two-tower models are attractive because item embeddings can often be precomputed offline and indexed for fast approximate nearest-neighbor retrieval.
At serving time, the system computes the user or session embedding and retrieves nearby items efficiently.

\subsection{Typical inputs}
The user/context tower may consume:
\begin{itemize}
  \item recent interaction history,
  \item device and channel context,
  \item customer segment,
  \item geographic or temporal features,
  \item active search or page context.
\end{itemize}

The item tower may consume:
\begin{itemize}
  \item item ID,
  \item category,
  \item brand,
  \item price bucket,
  \item title and description text,
  \item image-derived features,
  \item seller or inventory context.
\end{itemize}

\section{Approximate nearest-neighbor retrieval}
Deep retrieval only becomes operationally useful when the serving system can search the embedding space quickly.
This is where approximate nearest-neighbor (ANN) indexing enters the architecture.

\subsection{Why exact search is often too expensive}
If the platform must compare each request embedding against millions of item embeddings, exact retrieval may be too slow for production latency budgets.
ANN methods trade a small amount of exactness for major speed improvements.

\subsection{Architectural implication}
This means recommender design now depends on more than model training.
It also depends on:
\begin{itemize}
  \item embedding refresh cadence,
  \item index rebuild strategy,
  \item item insertion and deletion handling,
  \item consistency between item catalog state and retrieval index.
\end{itemize}

An item that is unavailable or unpublished should not remain in the active candidate index without explicit policy.

\section{Sequence models and session-aware recommendation}
Many commerce interactions are highly sequential.
A user may move from broad exploration to narrow intent within a single session.
Classical long-term preference models may miss this short-term drift.

\subsection{Session behavior}
Session-aware recommendation tries to capture recency and order, not only aggregate history.
Examples of useful sequential signals include:
\begin{itemize}
  \item last viewed products,
  \item order of category transitions,
  \item repeated brand examination,
  \item add-to-cart followed by substitution browsing,
  \item search reformulation patterns.
\end{itemize}

\subsection{Sequence model families}
Common sequence-oriented approaches include:
\begin{itemize}
  \item recurrent models,
  \item convolutional sequence models,
  \item Transformer-style attention models,
  \item next-item prediction objectives.
\end{itemize}

\subsection{Why sequence models help}
These models help because they can distinguish:
\begin{itemize}
  \item stable preference from temporary intent,
  \item exploration from purchase preparation,
  \item repeated exposure from genuine narrowing of interest.
\end{itemize}

This is especially useful for anonymous users, where long-term account history may be unavailable.

\section{Multi-modal recommendation}
Deep systems are particularly useful when products have rich unstructured content.
E-commerce catalogs often contain:
\begin{itemize}
  \item titles and descriptions,
  \item structured attributes,
  \item images,
  \item reviews,
  \item seller-generated content,
  \item taxonomy information.
\end{itemize}

\subsection{Text and semantic understanding}
Text encoders can help represent product meaning beyond exact attribute matching.
For example, two products may be semantically similar even when they do not share the same surface keywords.

\subsection{Image understanding}
Visual features matter strongly in apparel, furniture, decor, and other style-driven categories.
Image embeddings can support:
\begin{itemize}
  \item visually similar item recommendation,
  \item cold-start support for new products,
  \item quality or anomaly checks for catalog content.
\end{itemize}

\subsection{Fusion}
Modern item representations often combine ID-based, attribute-based, text-based, and image-based features.
This reduces dependence on raw interaction history alone and helps address sparse data problems.

\section{Deep re-ranking}
Retrieval is only the first stage.
Deep models can also be applied in re-ranking after candidates are retrieved.

\subsection{Why re-ranking differs from retrieval}
The re-ranker can use richer cross-features because it evaluates a much smaller candidate set.
For example, it may incorporate:
\begin{itemize}
  \item user-item interaction history,
  \item session context,
  \item merchandising rules,
  \item margin, inventory, or delivery constraints,
  \item diversity adjustments,
  \item fraud or return-risk features.
\end{itemize}

\subsection{Business role of re-ranking}
The re-ranker is often where business policy becomes explicit.
The retrieval stage finds promising candidates.
The re-ranker decides which candidates best satisfy a mix of relevance and operational constraints.

\section{Training data and objectives}
Deep recommender quality depends heavily on how the training problem is defined.

\subsection{Positive and negative examples}
Training often requires:
\begin{itemize}
  \item positive examples from clicks, carts, or purchases,
  \item sampled negatives from unseen or skipped items,
  \item sometimes hard negatives that were exposed but not chosen.
\end{itemize}

\subsection{Objective choices}
Common objectives include:
\begin{itemize}
  \item next-item prediction,
  \item pairwise ranking losses,
  \item contrastive retrieval objectives,
  \item multi-task objectives combining click, cart, and purchase behavior.
\end{itemize}

\subsection{Label quality}
As in earlier chapters, the label design is a business interpretation problem.
If a purchase is later returned, or if a click was caused mainly by position bias, the training data need careful handling.
Deep models can amplify weak signals just as efficiently as strong ones.

\section{Cold start in deep recommendation}
Deep systems do not eliminate cold start.
They only change the available tools for addressing it.

\subsection{User cold start}
For new users, deep models may rely on:
\begin{itemize}
  \item session behavior,
  \item page context,
  \item acquisition source,
  \item demographic or geo context where appropriate,
  \item popularity and trend priors.
\end{itemize}

\subsection{Item cold start}
For new items, deep models can use metadata and content features more effectively than classical ID-only approaches.
Text and image embeddings are especially helpful because they allow the item to be placed in a meaningful neighborhood before interaction data accumulate.

\subsection{Catalog churn}
Commerce catalogs change frequently.
A useful architecture must define how new item embeddings are generated, validated, indexed, and exposed without waiting for a full model retrain in every case.

\section{Diversity, exploration, and bias}
A recommender that always shows the most probable item is not necessarily the best commerce system.
The platform may need to balance relevance with diversity and exploration.

\subsection{Diversity}
Diversity can improve user experience by preventing repetitive lists and broadening discovery.
It may also improve business robustness by reducing over-concentration on a few items or sellers.

\subsection{Exploration}
Exploration is necessary to learn about new items, uncertain user interests, and shifting trends.
Without exploration, the platform over-commits to historical winners and limits future learning.

\subsection{Bias}
Deep systems can intensify:
\begin{itemize}
  \item popularity bias,
  \item seller-exposure bias,
  \item feedback-loop bias,
  \item and representation bias if product metadata are uneven or low quality.
\end{itemize}

This means governance remains essential even when model quality appears strong offline.

\section{Serving and latency constraints}
Deep recommenders exist under production latency budgets.
This makes serving architecture as important as modeling.

\subsection{Key serving questions}
\begin{itemize}
  \item Which features are available in real time?
  \item Which embeddings can be precomputed?
  \item How often are indexes refreshed?
  \item What is the fallback plan if the retrieval service is unavailable?
  \item How are out-of-stock or restricted items removed before final ranking?
\end{itemize}

\subsection{Caching and freshness}
Caching can reduce latency, but stale caches can recommend unavailable or irrelevant items.
Architecture must define freshness tolerances by surface.
For example, homepage suggestions may tolerate more staleness than cart recommendations near checkout.

\subsection{Fallback behavior}
Production systems usually need fallback strategies such as:
\begin{itemize}
  \item popularity-based lists,
  \item category best-sellers,
  \item recent-view heuristics,
  \item rule-based complements.
\end{itemize}

Fallbacks are not a sign of failure.
They are part of robust recommender system design.

\section{Evaluation beyond offline metrics}
Deep recommendation systems are often evaluated offline with ranking metrics, but their real value appears only in production.

\subsection{Why deep models can mislead offline}
Deep models may improve offline ranking scores because they capture rich structure, yet still disappoint online if:
\begin{itemize}
  \item latency increases and user experience degrades,
  \item popularity bias worsens,
  \item retrieval excludes important fresh inventory,
  \item serving features differ from training features,
  \item business constraints are ignored in ranking.
\end{itemize}

\subsection{Online considerations}
Relevant online metrics include:
\begin{itemize}
  \item click-through and add-to-cart lift,
  \item incremental conversion,
  \item revenue per session,
  \item item and seller coverage,
  \item return rate,
  \item latency and availability,
  \item fairness or exposure guardrails.
\end{itemize}

\section{A practical deep recommender stack}
A practical deep recommender architecture often includes:
\begin{itemize}
  \item event collection and identity stitching,
  \item a feature store or equivalent feature pipelines,
  \item model training jobs,
  \item embedding generation,
  \item ANN index building,
  \item retrieval serving,
  \item re-ranking serving,
  \item experimentation and monitoring layers.
\end{itemize}

This is why deep recommendation belongs not only to machine learning, but also to systems design and MLOps.

\section{Hands-on lab: retrieval plus re-ranking}
The goal of the practical work is to help students think beyond one-model recommendation.

\subsection{Lab scenario}
Students are given a dataset with user interaction histories, product metadata, and timestamps.
They are asked to design a two-stage pipeline for recommendation.

\subsection{Lab tasks}
\begin{enumerate}
  \item Train or simulate product embeddings from interaction and metadata signals.
  \item Build a simple candidate retrieval mechanism using embedding similarity.
  \item Add a lightweight re-ranking stage using business-aware features such as inventory or price sensitivity.
  \item Define a serving plan that explains which representations are precomputed and which features are computed online.
  \item Propose online guardrail metrics that would be monitored during an experiment.
\end{enumerate}

\subsection{Expected learning outcome}
By the end of the lab, students should be able to explain why a high-quality recommender is a pipeline and not just a neural model.

\section{Managerial implications}
Deep recommendation systems can unlock major gains in relevance and catalog understanding, but they also introduce higher system complexity.
Organizations should adopt them when the business problem justifies:
\begin{itemize}
  \item richer representation learning,
  \item very large-scale retrieval,
  \item multi-modal item understanding,
  \item session-aware adaptation.
\end{itemize}

They should not adopt them merely because they sound more advanced.
Without strong data foundations, business constraints, and serving discipline, deep recommenders can become expensive complexity with limited net value.

\section{Multiple-choice questions (MCQs)}
\begin{enumerate}
  \item In a modern recommender stack, the main purpose of \emph{candidate generation} is to:
  \begin{enumerate}
    \item finalize the exact business KPI improvement.
    \item reduce a very large catalog to a manageable subset before detailed ranking.
    \item eliminate the need for embeddings.
    \item replace experimentation.
  \end{enumerate}

  \item A two-tower retrieval model is most accurately described as:
  \begin{enumerate}
    \item a model with one tower for pricing and one tower for fraud.
    \item separate encoders for user/context and item representations whose similarity supports retrieval.
    \item a rule engine with two fallback policies.
    \item a batch-only ERP integration pattern.
  \end{enumerate}

  \item Why are embeddings useful in e-commerce recommendation?
  \begin{enumerate}
    \item They always eliminate cold start completely.
    \item They provide dense learned representations that support similarity and retrieval across users, items, and contexts.
    \item They make item metadata unnecessary.
    \item They are only relevant for payment systems.
  \end{enumerate}

  \item Sequence models are especially valuable when:
  \begin{enumerate}
    \item user intent changes within a session and order of interactions matters.
    \item the catalog has only ten products.
    \item recommendation is based only on ERP invoices.
    \item latency does not matter.
  \end{enumerate}

  \item Which is the best example of a production concern specific to deep retrieval systems?
  \begin{enumerate}
    \item Whether the catalog has chapter numbers.
    \item How to refresh embeddings and ANN indexes while respecting latency and availability constraints.
    \item Whether to use bullet points in documentation.
    \item Whether offline MAP@$K$ is equal to precision@$K$.
  \end{enumerate}
\end{enumerate}

\subsection*{Answer key}
\begin{enumerate}
  \item (b)
  \item (b)
  \item (b)
  \item (a)
  \item (b)
\end{enumerate}

\section{Exercises (short)}
\begin{enumerate}
  \item Compare a classical matrix-factorization recommender and a two-tower retrieval system for a large commerce catalog. State two reasons why the deep system may help and two costs it introduces.
  \item Propose an embedding feature set for items in a fashion store. Include structured, textual, and visual signals.
  \item Design a retrieval plus re-ranking pipeline for a cart recommendation widget. Specify where you would enforce inventory and diversity constraints.
\end{enumerate}

\section{Mini-case (Odoo-linked, vendor-neutral)}
\textbf{Scenario:} A hybrid B2C/B2B commerce company uses a headless storefront, Odoo for ERP, and a central ML platform. The product catalog is large and changes frequently.

\textbf{Task:} Design a deep recommender architecture:
\begin{itemize}
  \item Define the inputs to the user/context tower and the item tower.
  \item Explain how item embeddings would be refreshed when products, prices, or availability change.
  \item Identify two risks in using deep retrieval for B2B customers with negotiated catalogs or restricted assortments, and propose mitigations.
\end{itemize}
